\documentclass[11pt]{article}

% basic packages
\usepackage[margin=1in]{geometry}
\usepackage[pdftex]{graphicx}
\usepackage{amsmath,amssymb,amsthm}
\usepackage{tikz-cd}
\usepackage{william}

% page formatting
\usepackage{fancyhdr}
\pagestyle{fancy}

\renewcommand{\sectionmark}[1]{\markright{\textsf{\arabic{section}. #1}}}
\renewcommand{\subsectionmark}[1]{}
\lhead{\textbf{\thepage} \ \ \nouppercase{\rightmark}}
\chead{}
\rhead{}
\lfoot{}
\cfoot{}
\rfoot{}
\setlength{\headheight}{14pt}

\linespread{1.03}
\setlength{\parindent}{0.2in}
\setcounter{secnumdepth}{2}
\setcounter{tocdepth}{2}
\providecommand{\theHeexample}{\thesection.\arabic{eexample}}

\begin{document}

\thispagestyle{empty}
\bigskip \
\vspace{0.1cm}

\begin{center}
{\fontsize{22}{22} \selectfont Lecture Notes on}
\vskip 16pt
{\fontsize{36}{36} \selectfont \bf \sffamily Core Math}
\vskip 24pt
{\fontsize{18}{18} \selectfont \rmfamily \href{https://will-lancer.github.io}{Will Lancer}}
\vskip 6pt
{\fontsize{14}{14} \selectfont \ttfamily will.m.lancer@gmail.com}
\vskip 24pt
\end{center}

{\parindent0pt \baselineskip=15.5pt}
\noin
Lecture notes on core mathematics. This document now collects the algebra and
geometry/topology material in one place. Resources used:
\begin{itemize}
    \item Algebra:
    \begin{itemize}
        \item Aluffi's \emph{Algebra: Chapter 0}
        \item Atiyah and MacDonald's \emph{Introduction to Commutative Algebra}
    \end{itemize}
    \item Geometry and topology:
    \begin{itemize}
        \item Frederic Schuller's \emph{Geometrical Anatomy of Theoretical Physics}
        video lectures
        \item Bott and Tu's \emph{Differential Forms in Algebraic Topology}
        \item Littlejohn's \emph{Physics 250} lecture notes
    \end{itemize}
\end{itemize}

\newpage
\microtoc
\newpage

\section{Algebra}

\subsection{Categorical preliminaries}

\note{fill-in}

\subsection{Groups}

We review the properties of groups here. First, two high-brow definitions of a
group:
\begin{itemize}
    \item A \vocab{group} is a group object in \cat{Set}.
    \item A \vocab{group} is a one-object groupoid.
\end{itemize}
These aren't bad definitions, but they're more explicitly useful when they're
unpacked a bit.
\begin{definition}
    A \vocab{group} is a set with an associative binary operation
    $m_G \colon G \to G$ that has identites and inverses.
\end{definition}
We can phrase this by saying that a group has two additional unique maps,
$\iota \colon G \to G$ by $g \mapsto g^{-1}$ and $\id \colon G \to G$
by $g \mapsto g$. This is the more categorical and ``morally correct''
way of putting this.
\begin{eexample}
    [Canonical group examples]
    The most canonical group examples are
    \begin{itemize}
        \item Any field with addition, $\ints$, $\reals$, $\complex$, $\mathbb{Q}$,
        etc.. Some of them also work with multiplication.
        \item $\ints/n\ints$: the integers mod $n$ under addition. $(\ints/n\ints)^\times$
        are the units mod $n$.
        \item $D_n$. The isometries of a planar solid.
        \item $S_n$ and $A_n$. The permutations and even permutations of $n$ objects.
        \item Matrices over some field under addition or multiplication (if the operation
        is multiplication then needs invertibility).
        \item Types of functions under addition, e.g. smooth, periodic, etc..
    \end{itemize}
\end{eexample}

Groups form a category, \cat{Grp}. Its maps are \vocab{group homomorphisms},
which are maps $\varphi \colon G \to G'$ that respect the group structure,
$\varphi(ab) = \varphi(a)\varphi(b)$. One way of saying this is that the
following diagram commutes:\\
\note{put diagram}\\
We now investigate what other universal constructions exist in \cat{Grp}.
\begin{itemize}
    \item \textbf{Products}: the product in \cat{Grp} is the \vocab{direct product
    group}. This means that for all $\varphi_G \colon A \to G$, $\varphi_H \colon A \to H$,
    there exists a unique map $\widetilde{\varphi} \colon A \to G \times H$
    making the diagram
    \begin{center}
        \begin{tikzcd}
        & & G \\
        A \arrow[r, dashed, "\widetilde{\varphi}"] \arrow[urr, bend left=15, "\varphi_G"] \arrow[drr, bend right=15, "\varphi_H"'] & G \times H \arrow[ur, "\pi_G"] \arrow[dr, "\pi_H"'] & \\
        & & H
        \end{tikzcd}
    \end{center}
    commute. Concretely, taking the binary operations on $G$
    and $H$ and defining
    \begin{align*}
        m_G \times m_H \colon (G \times H) & \times (G \times H) \to G \times H\\
        (m_G \times m_H)((g_1, h_1), \, (g_2, h_2)) & \mapsto (m_G((g_1, g_2)), \, m_H((h_1, h_2))).
    \end{align*}
    gives the set $G \times H$ a group structure
    \begin{align*}
        (g_1, h_1) * (g_2, h_2) = (g_1 g_2, h_1 h_2).
    \end{align*}
    \item \textbf{Coproducts}: there is not a generic coproduct in \cat{Grp}.
    \note{talk about free product and \cat{Ab}}
    \item \textbf{Quotients} \note{add}
\end{itemize}

\section{Geometry and Topology}

\subsection{Manifolds}

\subsection{Differential constructions}

Define $\Omega^*$ to be the $\reals$-algebra generated by the symbols
$dx^1, \ldots, dx^n$ satisfying $dx^i dx^i = 0$ and
$dx^i dx^j = - dx^j dx^i$\footnote{This second condition is
actually necessary, as we're dealing with the multiplication of
the symbols themselves; we can't just apply condition (1)
to $dx^i + dx^j$ and get condition (2) for free.}. We define the
\vocab{vector space of $C^\infty$ differential forms} to be
\begin{align*}
    \Omega^*(M) = C^\infty(M) \otimes_\reals \Omega^*.
\end{align*}
A \vocab{differential $k$-form} is then a function in $\Omega^*(M)$
that has $k$ $dx^i$'s in it, i.e. $\omega = \omega_I dx^I$
for some multi-index $I = (i_1, \ldots, i_k)$. We can rewrite $\Omega^*(M)$
in terms of these spaces as $\Omega^*(M) = \bigoplus_{k = 0}^{\dim{M}} \Omega^k(M)$;
this is just saying that you can partition this space into subspaces
based on the degree of the differential form. Note the upper summand: by
anti-symmetry, $k \leq \dim{M}$ (because otherwise the forms would be
linearly dependent).

\begin{reemark}
    There is a clear similarity here between $\Omega^*(M)$ and the Fock
    space $\mathcal{F} = \bigoplus_{n = 0}^{\infty} \mathcal{H}_n$
    by analogizing $n$-forms to $n$-particle states. The similarity increases
    further when you consider a fermionic Fock space, where you directly
    identify fermionic fields $\psi \leftrightarrow dx$.
\end{reemark}

There are two natural constructions on the space of differential forms: the
\vocab{wedge product} and the \vocab{exterior derivative}. The wedge product is
just the natural way of multiplying forms; it's directly given by
\begin{align*}
    \boxed{\omega \wedge \eta = \omega_I \eta_J dx^I dx^J.}
\end{align*}
By definition, we can see that
$\omega \wedge \eta = (-1)^{\deg{\omega} \deg{\eta}} \, \eta \wedge \omega$,
as we pick up minus signs when we move the forms past each other.
The natural derivation on $\Omega^*(M)$ is the \vocab{exterior derivative},
$d \colon \Omega^n(M) \to \Omega^{n+1}(M)$, defined by
\begin{itemize}
    \item If $f \in \Omega^0(M)$, then $df = \partial_i f dx^i$.
    \item If $\omega \in \Omega^k(M)$, then $d\omega = df_I dx^I$.
\end{itemize}

The two most important properties of the exterior derivative are
that it's nilpotent and that it obeys graded Leibniz. That is,
\begin{align*}
    & d^2 \equiv 0
    & d(\omega \wedge \eta) = d\omega \wedge \eta + (-1)^{\deg{\omega}} \, \omega \wedge d\eta.
\end{align*}
Notice that the minus sign in the second term only depends on the degree
of the first form in the product. This is because you have to anti-commute
the $d$ only past $\omega$, thus moving it past $\deg{\omega}$ $dx$'s.

\begin{eexample}
    [Div, grad, curl, and all that]
    With the exterior derivative we can regain the definitions of
    div, grad, and curl extremely easily. On $\reals^3$, $\Omega^0(M)
    \cong C^\infty(\reals^3) \cong \Omega^3(M)$, and $\Omega^1(M) \cong \Omega^2(M)$,
    so we can identify these spaces with one another. Let $f \in \Omega^0(M)$.
    On functions,
    \begin{align*}
        df & = \partial_x f \, dx + \partial_y f \, dy + \partial_z f \, dz \leftrightarrow \grad f
    \end{align*}
    On one-forms,
    \begin{align*}
        d(f_i \, dx^i) & = (\partial_x f_y - \partial_y f_x) \, dx \wedge dy + (\partial_y f_z - \partial_z f_y) \, dy \wedge dz + (\partial_z f_x - \partial_x f_z) \, dz \wedge dx \leftrightarrow \curl f
    \end{align*}
    And on two-forms,
    \begin{align*}
        d(2\partial_{[i} f_{j]} \, dx^i \wedge dx^j) & = \left( \partial_x f_x + \partial_y f_y + \partial_z f_z \right) dx \wedge dy \wedge dz \leftrightarrow \div f
    \end{align*}
    In summary,
    \begin{align*}
        d(0\text{-forms}) & = \text{gradient},\\
        d(1\text{-forms}) & = \text{curl},\\
        d(2\text{-forms}) & = \text{divergence}.
    \end{align*}
\end{eexample}

Pretty cool!

\note{talk about cohomology theory here, then rewrite the first part of this
talking about how getting to cohomology is one of our explicit goals}

\end{document}
